\section*{अध्याय \ref{ch:g5:division} --- भाग और गुणज}

\begin{solution}{exo:g5:division:1}
$851 \div 23$: $23$ के \omterm{def:g5:division:multiple}{गुणज}: $23, 46, 69, 92, 115, 138, 161,
184, 207$। फिर $85 \div 23$: $3$ बार ($69$), \omterm{def:g3:division:remainder}{शेषफल} $16$;
$1$ नीचे उतारो: $161 \div 23 = 7$ ठीक-ठीक। \omterm{def:g3:division:remainder}{भागफल} $37$,
\omterm{def:g3:division:remainder}{शेषफल} $0$ ($23 \times 37 = 851$)।

$704 \div 32$: $70 \div 32$: $2$ बार ($64$), \omterm{def:g3:division:remainder}{शेषफल} $6$; $4$
नीचे उतारो: $64 \div 32 = 2$। \omterm{def:g3:division:remainder}{भागफल} $22$, \omterm{def:g3:division:remainder}{शेषफल} $0$।
\end{solution}

\begin{solution}{exo:g5:division:2}
$9 \div 2 = 4.5$; \quad $27 \div 4 = 6.75$; \quad
$33 \div 5 = 6.6$; \quad $21 \div 8 = 2.625$।
\end{solution}

\begin{solution}{exo:g5:division:3}
$54 \div 4 = 13.5$: हर एक $13.50$ देता है।
\end{solution}

\begin{solution}{exo:g5:division:4}
$10 \div 3 = 3.333\dots$: अंक $3$ हमेशा दोहराता रहता है ---
भाग कभी रुकता ही नहीं। सौवें भाग तक गोल: $3.33$।
\end{solution}

\begin{solution}{exo:g5:division:5}
$70$ तक $7$ के \omterm{def:g5:division:multiple}{गुणज}: $7, 14, 21, 28, 35, 42, 49, 56, 63,
70$। $18$ के \omterm{def:g5:division:multiple}{भाजक}: $1, 2, 3, 6, 9, 18$। $24$ के \omterm{def:g5:division:multiple}{भाजक}:
$1, 2, 3, 4, 6, 8, 12, 24$।
\end{solution}

\begin{solution}{exo:g5:division:6}
$470$: $2$ से (अंत में $0$), $5$ से, $10$ से; अंकों का योग $11$:
$3$ से नहीं।

$735$: अंत में $5$: $5$ से, $2$ या $10$ से नहीं; अंकों का योग
$15$: $3$ से।

$8\,001$: विषम; अंकों का योग $9$: केवल $3$ से।

$1\,314$: सम; अंकों का योग $9$: $2$ से और $3$ से (इसलिए $6$ से),
$5$ से नहीं।
\end{solution}

\begin{solution}{exo:g5:division:7}
$3$ और $5$ दोनों से विभाज्य यानी $15$ से विभाज्य: $60$ और $80$
के बीच संख्याएँ $60$, $75$।
\end{solution}

\begin{solution}{exo:g5:division:8}
$7.2 \div 6 = 1.2$: हर टुकड़ा $1.2$ m का है (जाँच:
$1.2 \times 6 = 7.2$)।
\end{solution}

\begin{solution}{exo:g5:division:9}
$1\,000 = 24 \times 41 + 16$: इकतालीस थैले पूरे भरते हैं, $16$
कंचे बचते हैं।
\end{solution}

\begin{solution}{exo:g5:division:10}
$7 \div 5 = 1.4$: पाँचों दोस्तों में हर एक को $1.4$ चॉकलेट
मिलती है --- चॉकलेट को दसवें भागों में काटकर यह हो सकता है।
$5 \div 7$ के लिए: भाग $0.714285\dots$ देता है, जो कभी रुकता
नहीं --- कोई ठीक-ठीक दशमलव हिस्सा है ही नहीं; केवल \omterm{def:g4:fractions:def}{भिन्न}
$\frac57$ उसे ठीक-ठीक नाम देती है।
\end{solution}

\begin{solution}{exo:g5:division:11}
अंकों का योग: $5 + 2 + d = 7 + d$। $3$ से विभाज्य तब, जब
$7 + d$, $9$, $12$ या $15$ हो: $d = 2$, $5$ या $8$। सबसे छोटा
$d = 2$ है ($522 = 3 \times 174$)।
\end{solution}
